\documentclass[12pt, a4paper]{article}
\usepackage[utf8]{inputenc}
\usepackage[ngerman]{babel}
\usepackage{amsmath, amssymb}
\usepackage{pgfplots}
\pgfplotsset{compat=1.18}
\usepackage{geometry}
\geometry{margin=2.5cm}
\usepackage{parskip}

\title{\textbf{Aufgabe 12 -- Lösung} \\ Methoden I: Mathematik, 2026S}
\author{}
\date{}

\begin{document}
\maketitle

\section*{Grundregel}
Der \textbf{Definitionsbereich} $D$ enthält alle $x$-Werte, für die die Funktion definiert ist. Probleme entstehen bei:
\begin{itemize}
    \item \textbf{Division durch 0} -- Nenner $\neq 0$
    \item \textbf{Wurzel aus negativer Zahl} -- Ausdruck unter $\sqrt{\phantom{x}} \geq 0$
\end{itemize}

% ============================================================
\section*{(a) $f_1(x) = \sqrt{\dfrac{x - 3}{(x-5)(x-1)}}$}

\textbf{Bedingung 1:} Nenner $\neq 0$: $(x-5)(x-1) \neq 0$, also $x \neq 5$ und $x \neq 1$.

\textbf{Bedingung 2:} Ausdruck unter der Wurzel $\geq 0$:
\[
    \frac{x - 3}{(x-5)(x-1)} \geq 0
\]

Nullstellen und Polstellen: $x = 1, 3, 5$. Diese teilen die Zahlengerade in Intervalle:

\begin{center}
\begin{tabular}{c|cccc}
    Intervall & $x - 3$ & $x - 5$ & $x - 1$ & Bruch \\
    \hline
    $x < 1$ & $-$ & $-$ & $-$ & $\frac{-}{(-)(-))} = \frac{-}{+} = -$ \\
    $1 < x < 3$ & $-$ & $-$ & $+$ & $\frac{-}{(-)(+)} = \frac{-}{-} = +$ \\
    $3 < x < 5$ & $+$ & $-$ & $+$ & $\frac{+}{(-)(+)} = \frac{+}{-} = -$ \\
    $x > 5$ & $+$ & $+$ & $+$ & $\frac{+}{(+)(+)} = +$
\end{tabular}
\end{center}

Der Bruch ist $\geq 0$ in den Intervallen $(1, 3]$ und $(5, \infty)$.

Bei $x = 3$ ist der Bruch $= 0$ (erlaubt: $\sqrt{0} = 0$).

Bei $x = 1$ und $x = 5$: Nenner $= 0$ (verboten).

\[
    \boxed{D = (1, 3] \cup (5, +\infty)}
\]

\begin{center}
\begin{tikzpicture}
\begin{axis}[
    axis lines=middle, xlabel={$x$}, ylabel={$y$},
    xmin=-1, xmax=9, ymin=-0.5, ymax=4,
    width=12cm, height=6cm, samples=200,
    restrict y to domain=0:4,
]
    \addplot[blue, thick, domain=1.01:3] {sqrt((x-3)/((x-5)*(x-1)))};
    \addplot[blue, thick, domain=5.01:9] {sqrt((x-3)/((x-5)*(x-1)))};
    \addplot[dashed, gray] coordinates {(1,-0.5) (1,4)};
    \addplot[dashed, gray] coordinates {(5,-0.5) (5,4)};
\end{axis}
\end{tikzpicture}
\end{center}

% ============================================================
\section*{(b) $f_2(x) = \sqrt{\dfrac{3x^2}{x^2 - 1}}$}

\textbf{Bedingung 1:} $x^2 - 1 \neq 0$, also $x \neq \pm 1$.

\textbf{Bedingung 2:} $\frac{3x^2}{x^2 - 1} \geq 0$

Der Zähler $3x^2 \geq 0$ immer (ist 0 nur bei $x = 0$).

Also muss der Nenner auch $\geq 0$ sein (oder Zähler $= 0$):
\begin{itemize}
    \item $x = 0$: Bruch $= 0$ \checkmark
    \item $x^2 - 1 > 0 \Leftrightarrow x^2 > 1 \Leftrightarrow |x| > 1 \Leftrightarrow x < -1$ oder $x > 1$
\end{itemize}

Bei $x = 0$ ist der Bruch $= 0$ (erlaubt), aber $0$ liegt im Intervall $(-1, 1)$, wo der Nenner negativ ist. Also: $\frac{0}{\text{negativ}} = 0 \geq 0$ -- das ist erlaubt!

Genauer: Für $-1 < x < 1$ (und $x \neq 0$) ist $x^2 - 1 < 0$ und $3x^2 > 0$, also Bruch $< 0$ -- nicht erlaubt.

\[
    \boxed{D = (-\infty, -1) \cup \{0\} \cup (1, +\infty)}
\]

\begin{center}
\begin{tikzpicture}
\begin{axis}[
    axis lines=middle, xlabel={$x$}, ylabel={$y$},
    xmin=-4, xmax=4, ymin=-0.5, ymax=5,
    width=12cm, height=6cm, samples=200,
    restrict y to domain=0:5,
]
    \addplot[blue, thick, domain=-4:-1.01] {sqrt(3*x^2/(x^2-1))};
    \addplot[blue, thick, domain=1.01:4] {sqrt(3*x^2/(x^2-1))};
    \addplot[only marks, mark=*, blue, mark size=2pt] coordinates {(0,0)};
    \addplot[dashed, gray] coordinates {(-1,-0.5) (-1,5)};
    \addplot[dashed, gray] coordinates {(1,-0.5) (1,5)};
\end{axis}
\end{tikzpicture}
\end{center}

% ============================================================
\section*{(c) $f_3(x) = \dfrac{7}{2x^3 - 5x^2 + 3x}$}

\textbf{Bedingung:} Nenner $\neq 0$. Faktorisieren:

\begin{align*}
    2x^3 - 5x^2 + 3x &= x(2x^2 - 5x + 3)
\end{align*}

Quadratische Formel für $2x^2 - 5x + 3 = 0$:
\[
    x = \frac{5 \pm \sqrt{25 - 24}}{4} = \frac{5 \pm 1}{4}
\]

Also $x = \frac{6}{4} = \frac{3}{2}$ oder $x = \frac{4}{4} = 1$.

Faktorisierung: $2x^3 - 5x^2 + 3x = x(x - 1)(2x - 3)$

Nenner $= 0$ bei $x = 0$, $x = 1$, $x = \frac{3}{2}$.

\[
    \boxed{D = \mathbb{R} \setminus \left\{0,\; 1,\; \frac{3}{2}\right\}}
\]

\begin{center}
\begin{tikzpicture}
\begin{axis}[
    axis lines=middle, xlabel={$x$}, ylabel={$y$},
    xmin=-2, xmax=3, ymin=-30, ymax=30,
    width=12cm, height=7cm, samples=300,
    restrict y to domain=-30:30,
]
    \addplot[blue, thick, domain=-2:-0.01] {7/(2*x^3 - 5*x^2 + 3*x)};
    \addplot[blue, thick, domain=0.01:0.99] {7/(2*x^3 - 5*x^2 + 3*x)};
    \addplot[blue, thick, domain=1.01:1.49] {7/(2*x^3 - 5*x^2 + 3*x)};
    \addplot[blue, thick, domain=1.51:3] {7/(2*x^3 - 5*x^2 + 3*x)};
    \addplot[dashed, gray] coordinates {(0,-30) (0,30)};
    \addplot[dashed, gray] coordinates {(1,-30) (1,30)};
    \addplot[dashed, gray] coordinates {(1.5,-30) (1.5,30)};
\end{axis}
\end{tikzpicture}
\end{center}

% ============================================================
\section*{(d) $f_4(x) = \sqrt{\dfrac{-3x^2 - 1}{7x + 3}}$}

\textbf{Bedingung 1:} $7x + 3 \neq 0$, also $x \neq -\frac{3}{7}$.

\textbf{Bedingung 2:} $\frac{-3x^2 - 1}{7x + 3} \geq 0$

\textbf{Zähler:} $-3x^2 - 1 = -(3x^2 + 1)$

Da $3x^2 + 1 \geq 1 > 0$ für alle $x$, ist der Zähler \textbf{immer negativ}: $-3x^2 - 1 < 0$.

Damit der Bruch $\geq 0$ wird, muss der Nenner auch negativ sein:

$7x + 3 < 0 \Leftrightarrow x < -\frac{3}{7}$

Aber: Der Bruch wird nie genau $= 0$ (Zähler ist nie 0), also nur strikt $> 0$ möglich.

\[
    \boxed{D = \left(-\infty,\; -\frac{3}{7}\right)}
\]

\begin{center}
\begin{tikzpicture}
\begin{axis}[
    axis lines=middle, xlabel={$x$}, ylabel={$y$},
    xmin=-4, xmax=1, ymin=-0.5, ymax=5,
    width=12cm, height=6cm, samples=200,
    restrict y to domain=0:5,
]
    \addplot[blue, thick, domain=-4:-0.44] {sqrt((-3*x^2 - 1)/(7*x + 3))};
    \addplot[dashed, gray] coordinates {(-3/7,-0.5) (-3/7,5)};
    \node[gray] at (axis cs:-0.43,4.5) {\small $x = -\frac{3}{7}$};
\end{axis}
\end{tikzpicture}
\end{center}

\end{document}
